\documentclass[11pt]{article}
\usepackage[a4paper,margin=26mm]{geometry}
\usepackage{amsmath,amssymb,mathtools}
\usepackage{unicode-math}
\setmainfont{Latin Modern Roman}
\setmathfont{Latin Modern Math}
\usepackage{xurl,hyperref}
\hypersetup{hidelinks,pdftitle={UnsolvedMath Problem 10400023}}
\newcommand{\sref}[2]{\hyperlink{source:#1:#2}{[S#2]}}
\newcommand{\cataloguescope}[1]{#1}
\begin{document}
% UnsolvedMath ID 10400023; source catalogue code AMR-103-0023
\setcounter{section}{10400022}
\section{Conjecture 2.1 — ([220, Problem 1.92 (N)]) Fd(ZK)/Fd+1(ZK) is torsion free for each d.}\label{10400023}
{\small\sffamily UnsolvedMath ID 10400023 · Topology}
\subsection{Definitions and mathematical statement}

\paragraph{Shared notation.}$\mathbb N=\{1,2,\ldots\}$, and $\mathbb Z,\mathbb Q,\mathbb R,\mathbb C$ denote the integers, rationals, reals and complex numbers. Any other conventions are specified locally.


A knot or link is a smooth oriented embedding of one or more circles in the oriented $3$-sphere, considered up to orientation-preserving ambient isotopy. Write $O$ for the unknot, which bounds an embedded disk. Let $R\mathcal K$ be the free $R$-module on oriented knot isotopy classes. Extend a knot invariant to singular knots with ordinary transverse double points by $v(K_\times)=v(K_+)-v(K_-)$. Let $F_d(R\mathcal K)$ be the submodule generated by the alternating resolution sums of knots with $d$ such points. An $R$-valued finite-type (Vassiliev) invariant of degree at most $d$ is an $R$-linear map $R\mathcal K\to R$ vanishing on $F_{d+1}(R\mathcal K)$.
\paragraph{Statement recorded in the supplied source.}
For every integer $d\ge0$, is the abelian group
\[F_d(\mathbb Z\mathcal K)/F_{d+1}(\mathbb Z\mathcal K)\]
torsion-free? This means that $m x=0$ with integer $m\ne0$ forces $x=0$. The statement records the source's conjecture; no present-day status claim is inferred from its historical formulation. \sref{10400023}{2}
\subsection{Short English statement}
Does each graded piece of the integral finite-type knot filtration have no nonzero torsion elements?
\subsection{Sources}
\begin{description}
\item[\hypertarget{source:10400023:1}{S1}] ULAM. AI and the UnsolvedMath Contributors, \emph{UnsolvedMath: A Curated Collection of Open Mathematics Problems}, record 10400023 (AMR-103-0023). CC BY 4.0. \url{https://huggingface.co/datasets/ulamai/UnsolvedMath}.
\item[\hypertarget{source:10400023:2}{S2}] T. Ohtsuki (editor), \emph{Problems on invariants of knots and 3-manifolds}, Geometry \& Topology Monographs 4 (2002), pp. 377--572 (published 2004), item 2.1 and its surrounding definitions and remarks. \url{https://arxiv.org/abs/math/0406190v1}.
\end{description}
\label{end:10400023}
\end{document}
